\documentclass[12pt,letterpaper]{article}
\usepackage{graphicx,textcomp}
\usepackage{natbib}
\usepackage{setspace}
\usepackage{fullpage}
\usepackage{color}
\usepackage[reqno]{amsmath}
\usepackage{amsthm}
\usepackage{fancyvrb}
\usepackage{amssymb,enumerate}
\usepackage[all]{xy}
\usepackage{endnotes}
\usepackage{lscape}
\newtheorem{com}{Comment}
\usepackage{float}
\usepackage{hyperref}
\newtheorem{lem} {Lemma}
\newtheorem{prop}{Proposition}
\newtheorem{thm}{Theorem}
\newtheorem{defn}{Definition}
\newtheorem{cor}{Corollary}
\newtheorem{obs}{Observation}
\usepackage[compact]{titlesec}
\usepackage{dcolumn}
\usepackage{tikz}
\usetikzlibrary{arrows}
\usepackage{multirow}
\usepackage{xcolor}
\newcolumntype{.}{D{.}{.}{-1}}
\newcolumntype{d}[1]{D{.}{.}{#1}}
\definecolor{light-gray}{gray}{0.65}
\usepackage{url}
\usepackage{listings}
\usepackage{color}
\usepackage{relsize}

\definecolor{codegreen}{rgb}{0,0.6,0}
\definecolor{codegray}{rgb}{0.5,0.5,0.5}
\definecolor{codepurple}{rgb}{0.58,0,0.82}
\definecolor{backcolour}{rgb}{0.95,0.95,0.92}

\lstdefinestyle{mystyle}{
	backgroundcolor=\color{backcolour},   
	commentstyle=\color{codegreen},
	keywordstyle=\color{magenta},
	numberstyle=\tiny\color{codegray},
	stringstyle=\color{codepurple},
	basicstyle=\footnotesize,
	breakatwhitespace=false,         
	breaklines=true,                 
	captionpos=b,                    
	keepspaces=true,                 
	numbers=left,                    
	numbersep=5pt,                  
	showspaces=false,                
	showstringspaces=false,
	showtabs=false,                  
	tabsize=2
}
\lstset{style=mystyle}
\newcommand{\Sref}[1]{Section~\ref{#1}}
\newtheorem{hyp}{Hypothesis}


\title{Problem Set 4}
\date{Due: November 27, 2025}
\author{Rosalie LaCroix}


\begin{document}
	\maketitle
	\section*{Instructions}
	\begin{itemize}
		\item Please show your work! You may lose points by simply writing in the answer. If the problem requires you to execute commands in \texttt{R}, please include the code you used to get your answers. Please also include the \texttt{.R} file that contains your code. If you are not sure if work needs to be shown for a particular problem, please ask.
		\item Your homework should be submitted electronically on GitHub.
		\item This problem set is due before 23:59 on Thursday November 27, 2025. No late assignments will be accepted.
	\end{itemize}



	\vspace{.5cm}
\section*{Question 1: Economics}
\vspace{.25cm}
\noindent 	
In this question, use the \texttt{prestige} dataset in the \texttt{car} library. First, run the following commands:

\begin{verbatim}
install.packages(car)
library(car)
data(Prestige)
help(Prestige)
\end{verbatim} 


\noindent We would like to study whether individuals with higher levels of income have more prestigious jobs. Moreover, we would like to study whether professionals have more prestigious jobs than blue and white collar workers.

\begin{enumerate}
	
	\item [(a)]
	Create a new variable \texttt{professional} by recoding the variable \texttt{type} so that professionals are coded as $1$, and blue and white collar workers are coded as $0$ (Hint: \texttt{ifelse}).
	
	\lstinputlisting[language=R, firstline=43, lastline=44]{PS04_RL.R}
	
	\vspace{2cm}
	
	\item [(b)]
	Run a linear model with \texttt{prestige} as an outcome and \texttt{income}, \texttt{professional}, and the interaction of the two as predictors (Note: this is a continuous $\times$ dummy interaction.)
	
	\lstinputlisting[language=R, firstline=48, lastline=49]{PS04_RL.R}
	
	\textsmaller[1]{There is a positive, significant linear relationship between income and prestige score, as well as between type of occupation (Professional) and prestige score. The Professional variable is a dummy variable defined as blue and white collar workers = 0, and professional, managerial, and technical workers = 1. There is a negative, statistically significant linear relationship between the interaction variable income*professional and prestige score. This means that when professional is 1, the impact of income on prestige score decreases by 0.002.}
	
	\begin{table}[!htbp] \centering 
		\caption{Table of Coefficients} 
		\label{} 
		\begin{tabular}{@{\extracolsep{5pt}}lc} 
			\\[-1.8ex]\hline 
			\hline \\[-1.8ex] 
			& \multicolumn{1}{c}{\textit{Dependent variable:}} \\ 
			\cline{2-2} 
			\\[-1.8ex] & Prestige \\ 
			& Model 1 \\ 
			\hline \\[-1.8ex] 
			Income & 0.003$^{***}$ \\ 
			& (0.0005) \\ 
			& \\ 
			Professional & 37.781$^{***}$ \\ 
			& (4.248) \\ 
			& \\ 
			Income:Professional & $-$0.002$^{***}$ \\ 
			& (0.001) \\ 
			& \\ 
			Constant & 21.142$^{***}$ \\ 
			& (2.804) \\ 
			& \\ 
			\hline \\[-1.8ex] 
			Observations & 98 \\ 
			R$^{2}$ & 0.787 \\ 
			Adjusted R$^{2}$ & 0.780 \\ 
			Residual Std. Error & 8.012 (df = 94) \\ 
			F Statistic & 115.878$^{***}$ (df = 3; 94) \\ 
			\hline 
			\hline \\[-1.8ex] 
			\textit{Note:}  & \multicolumn{1}{r}{$^{*}$p$<$0.1; $^{**}$p$<$0.05; $^{***}$p$<$0.01} \\ 
		\end{tabular} 
	\end{table} 
	
	\vspace{6cm}
	\item [(c)]
	Write the prediction equation based on the result.
	
		\textsmaller[1]{The prediction equation based on this result is y = 21.142 + 0.003$x_{1}$ + 37.781$D_{1}$ - 0.002$D_{1}$$x_{1}$. More specifically, for this particular model, prestige = 21.142 + 0.003*income + 37.781*professional - 0.002*income:professional.}
	\vspace{2cm}
	\item [(d)]
	Interpret the coefficient for \texttt{income}.
	
		\textsmaller[1]{The coefficient for income indicates that, on average, when professional is 0, for every one unit increase in income, there is a 0.003 increase in prestige. This can also be interpreted as, on average, for blue and white collar workers, for every one unit (\$1) increase in income, there is a 0.003 increase in Pineo-Porter prestige score for occupation.}
	
	\vspace{2cm}	
	\item [(e)]
	Interpret the coefficient for \texttt{professional}.
	
		\textsmaller[1]{The coefficient for professional indicates that, on average, when income is 0, for a one unit increase in professional, there is a 37.781 increase in prestige. This can also be interpreted as, on average, when income is 0, moving from a blue or white collar worker to a professional, results in a 37.781 increase in Pineo-Porter prestige score for occupation.}
	
	\newpage
	\item [(f)]
	What is the effect of a \$1,000 increase in income on prestige score for professional occupations? In other words, we are interested in the marginal effect of income when the variable \texttt{professional} takes the value of $1$. Calculate the change in $\hat{y}$ associated with a \$1,000 increase in income based on your answer for (c).
	
	\textsmaller[1]{Prediction equation from (c): prestige = 21.142 + 0.003*income + 37.781*professional - 0.002*income:professional.}

	\textsmaller[1]{Income\_0 = \$0; Income\_1000 = \$1,000; Professional = 1}
	
	\lstinputlisting[language=R, firstline=64, lastline=66]{PS04_RL.R}
	
	\textsmaller[1]{The marginal effect of a \$1,000 increase in income on prestige score for professional occupations is 1. This is an incredibly small effect on prestige score due to the small coefficient associated with a change in income, which is made even smaller due to the 0.002 decrease in income's effect on prestige score, caused by it's interaction with professional status when professional takes the value of 1.}
	
	\vspace{2cm}
	
	\item [(g)]
	What is the effect of changing one's occupations from non-professional to professional when her income is \$6,000? We are interested in the marginal effect of professional jobs when the variable \texttt{income} takes the value of $6,000$. Calculate the change in $\hat{y}$ based on your answer for (c).
	
		\textsmaller[1]{Prediction equation from (c): prestige = 21.142 + 0.003*income + 37.781*professional - 0.002*income:professional.}
	
	\textsmaller[1]{Non-prof: Income = \$6,000; Professional = 0}
	\textsmaller[1]{Prof: Income = \$6,000; Professional = 1}
	
	\lstinputlisting[language=R, firstline=72, lastline=74]{PS04_RL.R}
	
	\textsmaller[1]{The effect of changing one's occupation from non-professional to professional with an income of \$6,000 is a 25.781 increase in prestige score. Compared to a change in income, a change in professional status has a larger effect on prestige score due to the larger average increase in prestige score associated with the change from non-professional to professional jobs. It is important to note that for professional jobs, there is a decrease in the positive effect of income on prestige score, which resulted in the subtraction of 12 points when the income value is \$6,000. }
	
	
\end{enumerate}

\newpage

\section*{Question 2: Political Science}
\vspace{.25cm}
\noindent 	Researchers are interested in learning the effect of all of those yard signs on voting preferences.\footnote{Donald P. Green, Jonathan	S. Krasno, Alexander Coppock, Benjamin D. Farrer,	Brandon Lenoir, Joshua N. Zingher. 2016. ``The effects of lawn signs on vote outcomes: Results from four randomized field experiments.'' Electoral Studies 41: 143-150. } Working with a campaign in Fairfax County, Virginia, 131 precincts were randomly divided into a treatment and control group. In 30 precincts, signs were posted around the precinct that read, ``For Sale: Terry McAuliffe. Don't Sellout Virgina on November 5.'' \\

Below is the result of a regression with two variables and a constant.  The dependent variable is the proportion of the vote that went to McAuliff's opponent Ken Cuccinelli. The first variable indicates whether a precinct was randomly assigned to have the sign against McAuliffe posted. The second variable indicates
a precinct that was adjacent to a precinct in the treatment group (since people in those precincts might be exposed to the signs).  \\

\vspace{.5cm}
\begin{table}[!htbp]
	\centering 
	\textbf{Impact of lawn signs on vote share}\\
	\begin{tabular}{@{\extracolsep{5pt}}lccc} 
		\\[-1.8ex] 
		\hline \\[-1.8ex]
		Precinct assigned lawn signs  (n=30)  & 0.042\\
		& (0.016) \\
		Precinct adjacent to lawn signs (n=76) & 0.042 \\
		&  (0.013) \\
		Constant  & 0.302\\
		& (0.011)
		\\
		\hline \\
	\end{tabular}\\
	\footnotesize{\textit{Notes:} $R^2$=0.094, N=131}
\end{table}

\vspace{.5cm}
\begin{enumerate}
	
	\newpage
	
	\item [(a)] Use the results from a linear regression to determine whether having these yard signs in a precinct affects vote share (e.g., conduct a hypothesis test with $\alpha = .05$).
	
		\textsmaller[1]{I will utilize a t-test to test the partial effect of having yard signs in a precinct on vote share.}
	
		\textsmaller[1]{$H_{0}$ : $\beta_{1}$ = 0, $H_{a}$ : $\beta_{1}$ $\neq$ 0.}
	
		\lstinputlisting[language=R, firstline=80, lastline=88]{PS04_RL.R}
	
		\textsmaller[1]{We can reject the null hypothesis that having yard signs in a precinct does not effect vote share, meaning that the presence of yard signs does effect vote share. We can conclude this due to the fact that the p-value is 0.014, which is smaller than the alpha value of 0.05, making the effect of 0.042 on vote share significant.}
	
	\vspace{2cm}
	
	\item [(b)]  Use the results to determine whether being
	next to precincts with these yard signs affects vote
	share (e.g., conduct a hypothesis test with $\alpha = .05$).
	
		\textsmaller[1]{I will utilize a t-test to test the partial effect of being next to a precinct with yard signs on vote share.}
	
		\textsmaller[1]{$H_{0}$ : $\beta_{2}$ = 0, $H_{a}$ : $\beta_{2}$ $\neq$ 0.}
		
		\lstinputlisting[language=R, firstline=92, lastline=100]{PS04_RL.R}
		
		\textsmaller[1]{We can reject the null hypothesis that being next to precincts with yard signs does not effect vote share, meaning that the presence of yard signs does effect vote share. We can conclude this due to the fact that the p-value is 0.002, which is smaller than the alpha value of 0.05, making the effect of 0.042 on vote share significant.}
	
	\newpage
	
	\item [(c)] Interpret the coefficient for the constant term substantively.
	
		\textsmaller[1]{When there are no lawn signs present in treatment group precincts, and therefore no exposure to lawn signs for people in adjacent precincts, the vote share will be equal to 0.302.}
	
	\vspace{2cm}
	
	\item [(d)] Evaluate the model fit for this regression.  What does this	tell us about the importance of yard signs versus other factors that are not modeled?
	
		\textsmaller[1]{I will utilize a global F-test to evaluate whether or not the model containing the effects of yard signs in and next to a precinct is useful in explaining variation in vote share. This will test the full model, containing the two explanatory variables for presence of yard signs, against a reduced model containing only the constant.}
		
		\textsmaller[1]{Reduced model: Y = $\beta_{0}$ + $\epsilon$}
		
		\textsmaller[1]{Full model: Y = $\beta_{0}$ + $\beta_{1}$$X_{1}$ + $\beta_{2}$$X_{2}$ + $\epsilon$}
	
		\textsmaller[1]{$H_{0}$: $\beta_{1}$ = $\beta_{2}$ = 0, $H_{a}$ : At least one $\beta$ $\neq$ 0.}
	
		\lstinputlisting[language=R, firstline=104, lastline=114]{PS04_RL.R}
	
		\textsmaller[1]{The p-value for the F-test evaluating this model was found to be 0.0018, with a test statistic of 6.64, which indicates that we can reject the null hypothesis that both $\beta_{1}$ and $\beta_{2}$ are equal to 0. Rejecting the null hypothesis means that either the presence of yard signs in a precinct and/or in an adjacent precinct is not 0 and is therefore significant in predicting vote share. This means that other factors that are not modeled may not be more important in predicting vote share than the presence of yard signs and that the proposed full model is sufficient in predicting vote share.}
	
\end{enumerate}  


\end{document}
